% Generated draft: requires mathematical and source review.
\documentclass[11pt]{article}
\usepackage[margin=25mm]{geometry}
\usepackage{fontspec}
\setmainfont{Latin Modern Roman}
\usepackage{amsmath,amssymb,mathtools,unicode-math}
\setmathfont{Latin Modern Math}
\usepackage{enumitem,needspace,xurl,hyperref}
\setlength{\parindent}{0pt}
\setlength{\parskip}{5pt}
\setlength{\emergencystretch}{3em}
\newcommand{\N}{\mathbb{N}}
\newcommand{\Z}{\mathbb{Z}}
\newcommand{\Q}{\mathbb{Q}}
\newcommand{\R}{\mathbb{R}}
\newcommand{\C}{\mathbb{C}}
\newcommand{\F}{\mathbb{F}}
\newcommand{\A}{\mathbb{A}}
\newcommand{\PP}{\mathbb{P}}
\newcommand{\E}{\mathbb{E}}
\newcommand{\eps}{\varepsilon}
\newcommand{\dd}{\,\mathrm{d}}
\begin{document}
\section*{UnsolvedMath ID: 20003163}
\textbf{Source problem code:} catalogue entry
\subsection{Definitions and mathematical statement}
For the intended nontrivial isometric action of a compact Lie group $G$ on a Riemannian manifold $M$, polar means that a connected complete immersed section meets every orbit orthogonally. Cohomogeneity is the dimension of the orbit space for that action. The source asks whether a polar $M$ with strictly positive sectional curvature must be diffeomorphic to a compact rank-one symmetric space (CROSS). The excerpt does not specify compactness, simple connectivity, or the action's cohomogeneity. Trivial actions must be excluded in a meaningful rigidity formulation; otherwise every manifold is polar by taking itself as section.

The aid supplies a closed simply connected counterexample at cohomogeneity one: the Grove--Verdiani--Ziller manifold $P_2$, a $2$-connected smooth seven-manifold with a strictly positively curved $SO(4)$-invariant cohomogeneity-one metric. A cohomogeneity-one isometric action has a normal-geodesic section and is polar for that same metric. The reported invariant
\[
\pi_3(P_2)\cong\Z/2\Z
\]
excludes the seven-dimensional CROSS possibilities: $S^7$, and $\R\mathrm P^7$ if non-simply-connected rank-one spaces are included, both have zero third homotopy group. Thus the literal question is false even after adding closedness and simple connectivity.

The supplied repaired theorem is attributed to Fang--Grove--Thorbergsson: a simply connected closed positively curved manifold carrying a nontrivial polar action of cohomogeneity at least two, in that theorem's precise action hypotheses, is diffeomorphic to a CROSS. A complete theorem statement and citation are not present.

The independent simple-connectivity test is explicit. For integers $m\geq3$ and odd $p\geq3$, let
\[
M=L^{2m-1}(p;1,\ldots,1)=S^{2m-1}/\mu_p,
\]
where $\mu_p$ acts by simultaneous scalar multiplication on $\C^m$. The round metric descends with sectional curvature one. The coordinate phase action descends to the effective torus $T^m/\mu_p$. Its principal orbits have dimension $m$, hence cohomogeneity $m-1\geq2$. The image of the real unit sphere $S^{m-1}\subset\R^m$ is a complete section: coordinate phases move every orbit to it, and its real tangent directions are orthogonal to the imaginary orbit directions. Odd $p$ makes its projection injective. Yet $\pi_1(M)=\Z/p\Z$, so it is not a CROSS, whose relevant fundamental groups are trivial or order two. This preserves cohomogeneity at least two while violating simple connectivity.

Accordingly $P_2$ and the lens family separately show why the two added rigidity hypotheses matter. They do not replace the need to state the action and exact scope of the attributed theorem, nor classify every positively curved polar action.
\subsection{Short English statement}
The supplied examples refute the unqualified question and separate two missing hypotheses. The simply connected $P_2$ example has cohomogeneity one, while the round lens examples have cohomogeneity at least two and nontrivial fundamental group. The repaired simply connected higher-cohomogeneity rigidity theorem is attributed to Fang--Grove--Thorbergsson, with its full action hypotheses still requiring citation review.
\subsection{Sources}
\begin{itemize}
\item Local registry: permanent catalogue entry, exact source code catalogue entry; \texttt{aim-topology-notes.json}, \texttt{notes[250]}.
\item AIM workshop \emph{Manifolds with nonnegative sectional curvature}, Problem 38, \url{https://aimath.org/WWN/nnsectcurvature/nnsectcurvature.pdf}. No proposer, year, remarks, or original literature reference is supplied.
\item Supplied attributions: Grove--Verdiani--Ziller for the positively curved $P_2$ example and Fang--Grove--Thorbergsson for the repaired higher-cohomogeneity rigidity theorem. No titles, theorem numbers, or complete authoritative citations are included.
\item Supplied research attempt 1, registry updated 14 August 2026: counterexample, refuted-at-run label, low novelty confidence, machine-generated and not peer reviewed. It supplies the lens-space sharpness family and reports the $P_2$ homotopy invariant. Verification identifies Problem 38 but provides no full action conventions.
\end{itemize}
\end{document}
